\section{A precise reduction of the remaining dimensional bound}
\label{sec:highdimboundreduction}
The construction above and the contrast bound leave a specific structural
question. The next proposition isolates a sufficient additional inequality;
it does not assert that all the original $F$ tests force that inequality
in arbitrary dimension.

\begin{proposition}[Conditional three-state reduction]
\label{prop:highdimboundreduction}
Let $n\geq4$, let the energies be distinct and ordered, and let a positive
reversible generator at $0<\beta_b<\infty$ have a simple slow mode.
Suppose every original $G_{ijk}\leq0$. Orient the slow left eigenfunction
$w$ by its nonnegative weighted-prefix contrast coordinates, as in
Theorem~\ref{thm:orderedk}. If additionally
\begin{equation}\label{eq:highdimarithmetic}
w_1+w_2\geq2w_3,
\end{equation}
then every finite cooling pair $0<\beta_h<\beta_c<\beta_b$ satisfies
\[
\frac{|b_2(\beta_c)|}{|b_2(\beta_h)|}<\frac n3.
\]
In particular its amplitude advantage is strictly below $1-3/n$.
\end{proposition}
\begin{proof}
Write $q_i=\pi_i(\beta_b)$, $Q_k=\sum_{i\leq k}q_i$, and
\[
h^{(k)}_i=
\begin{cases}
1,&i\leq k,\\
-Q_k/q_{k+1},&i=k+1,\\
0,&i>k+1.
\end{cases}
\]
The contrast theorem supplies $w=\sum_{k=1}^{n-1}x_kh^{(k)}$ with
$x_k\geq0$, not all zero. Put
$f=x_1h^{(1)}+x_2h^{(2)}$; it is supported on the first three states.
Every contrast with $k\geq3$ is constant on these states, so
\[
f_1-f_2=w_1-w_2=\frac{Q_2}{q_2}x_1\geq0,
\qquad f_1+f_2-2f_3=w_1+w_2-2w_3\geq0.
\]
Thus the first three entries of $f$ are a constant plus
$u(1,-1,0)+v(1,1,-2)$ with $u,v\geq0$.

Let $m_f(\beta)$ be the conditional Gibbs mean of $f$ on states
$1,2,3$. This mean is increasing when $f\ne0$. Indeed, shifting
$E_1=0$ and writing $a=E_2$, $b=E_3$, $t=e^{-a\beta}$,
$y=e^{-b\beta}$, the conditional probabilities are
$(p_1,p_2,p_3)=(1,t,y)/(1+t+y)$ and
\[
(p_1-p_2)'=
\frac{2at+by-(b-a)ty}{(1+t+y)^2}>0,
\qquad (p_1+p_2-2p_3)'=-3p_3'>0.
\]
Here $0<a<b$ and $0<t\leq1$ for $\beta\geq0$.
If $f\ne0$, its zero bath mean prevents its first three entries from
being constant, so at least one of $u,v$ is positive. Consequently
$m_f(\beta_b)=0$ and $m_f(\beta)<0$ for $0\leq\beta<\beta_b$.

For $j\leq n$, write $P_j(\beta)=Z_j(\beta)/Z_n(\beta)$ for the
Gibbs mass on the first $j$ states. If $j<n$, their lower energies give
$j/n<P_j(\beta)<1$ for finite $\beta>0$. The overlap of $f$ is
$P_3m_f$, and hence, when $f\ne0$,
\[
0<\frac{|\langle f\rangle_{\beta_c}|}
        {|\langle f\rangle_{\beta_h}|}
=\frac{P_3(\beta_c)}{P_3(\beta_h)}
  \frac{|m_f(\beta_c)|}{|m_f(\beta_h)|}
<\frac n3.
\]

For each $k\geq3$, the conditional mean of $h^{(k)}$ on its first
$k+1$ states equals
\[
m_k(\beta)=1-\left(1+\frac{Q_k}{q_{k+1}}\right)
 \frac{e^{-\beta E_{k+1}}}{Z_{k+1}(\beta)}.
\]
It is strictly increasing, since $E_{k+1}$ is the largest energy in
that conditional system, and it vanishes at $\beta_b$. Its cooling
values are therefore negative. The same factorization gives
\[
\frac{|\langle h^{(k)}\rangle_{\beta_c}|}
     {|\langle h^{(k)}\rangle_{\beta_h}|}
<\frac{n}{k+1}\leq\frac n4.
\]
For $k+1=n$, the mass factor is identically one and strictness instead
comes directly from the strict increase of the negative conditional
mean; this handles the full-support contrast without asserting
$P_n<1$.

The nonzero terms in
$b_2(\beta)=\langle f\rangle_\beta+
\sum_{k=3}^{n-1}x_k\langle h^{(k)}\rangle_\beta$
are all negative on cooling. Its ratio of absolute overlaps is thus a
weighted average of the corresponding component ratios, with positive
weights proportional to their hot absolute overlaps. Each ratio is
strictly below $n/3$, proving the assertion.
\end{proof}

To state the remaining gap precisely, let $S_n$ be the supremum of
\[
1-\frac{|b_2(\beta_h)|}{|b_2(\beta_c)|}
\]
over positive reversible $n$-state generators with distinct ordered
energies, a positive finite bath, a simple slow mode, all original
$F_{ijk},G_{ijk}\leq0$, and finite cooling pairs. The construction in
Theorem~\ref{thm:highdimamplitude} and the general contrast bound give
\[
1-\frac3n\ \leq\ S_n\ \leq\ 1-\frac2n.
\]
Every individual finite pair has advantage strictly below the displayed
upper bound; that strictness does not replace the non-strict inequality
for the supremum. For $n=4$, Theorem~\ref{thm:fulltripletbound} settles
$S_4=1/4$. For $n>4$, proving that the full original triplet premise
forces \eqref{eq:highdimarithmetic} would, by the proposition and the
existing construction, settle $S_n=1-3/n$. This all-dimensional
arithmetic implication is not proved here. A counterexample to that
intermediate implication alone would not establish that the amplitude
supremum exceeds $1-3/n$; a thermal-overlap violation would still be needed.

\begin{proposition}[A constant rare-tail subfamily]
Let a positive reversible $n$-state generator, $n\geq4$, have distinct
ordered energies, a positive finite bath, and all original $F,G\leq0$.
Suppose its oriented simple slow eigenfunction is constant on states $4,\ldots,n$, and
$Q_T:=\sum_{k=4}^n q_k<q_3$. Then $w_1+w_2\geq2w_3$ and the
amplitude bound of Proposition~\ref{prop:highdimboundreduction} holds.
The full dynamics need not be lumpable.
\end{proposition}
\begin{proof}
Orthogonally compress the reversible backward generator onto functions
constant on the tail, using its $q$-weighted inner product. The compressed
four-state generator has weights $(q_1,q_2,q_3,Q_T)$, unchanged low-low
rates, and
\[
a_{iT}=\sum_{k\geq4}\frac{q_k}{Q_T}a_{ik},\qquad
a_{Ti}=\sum_{k\geq4}a_{ki}=\frac{Q_T}{q_i}a_{iT}.
\]
Its rates are positive, its diagonal conserves probability, and it is
reversible. The original slow eigenfunction belongs to this subspace,
so it remains an eigenfunction. The Rayleigh variational principle bounds
every compressed nonstationary eigenvalue above by the original slow
eigenvalue. Equality is simple because the original slow eigenspace is
one-dimensional. Hence this is also the compressed simple slow mode.

For $i<j\leq3$, direct substitution gives
\[
G_{ijT}=\sum_{k\geq4}\frac{q_k}{Q_T}G_{ijk},\qquad
F_{ijT}=\sum_{k\geq4}\frac{q_k}{Q_T}F_{ijk}
 +\sum_{k\geq4}\frac{q_k(Q_T-q_k)}{2Q_Tq_j}G_{ijk}.
\]
All coefficients are nonnegative, and the test on $123$ is unchanged.
The compressed weights are strictly ordered because $Q_T<q_3$.
The four-state arithmetic lemma therefore applies. Its orientation agrees
with the original one: the first two contrast coordinates are unchanged,
and the third is nonnegative because the common tail value is $w_n\leq0$.
The preceding conditional bound completes the proof.
\end{proof}

This compression is different from eliminating states by a Schur
complement. Its test preservation is the displayed identity, not an
assumption about effective rates. The restrictions of a constant tail
mode and $Q_T<q_3$ are substantive; aggregation can reorder energies
when the latter fails.
